% !PRINTED 602
\weberpage{602}

The constant $g$ is more closely determined by (5):

\weq{7}{g = \frac{\gamma\,\pi}{M^2\sqrt{-\Delta}},}

where $\gamma$ is the number of classes of numbers contained in $A$ with respect to the modulus $M$.

If by $T$ we wish to understand the number of integral ideals of the principal class $\overline{E}A$ whose norm is not greater than $t$, then (6) still holds in that case too; only, if $-1$ occurs in $A$, the number $g$ has to be halved, and in the two exceptional cases $\Delta = -4$, $\Delta = -3$ it may possibly have to be divided further by 4 or by 6.

\begin{center}
\textbf{\S\,163.\qquad \gesperrt{The Dirichlet Sums.}}
\end{center}

We now understand by $\overline{O}$ the ideal group of all the ideals in $\Omega$ (with regard to the excludent $S$); let $A$ be a congruence group and

\weq{1}{\overline{A} = \overline{E}A,}

the group of the corresponding principal ideals. We decompose $\overline{O}$ according to \S\,161, (3) into the \gesperrt{classes} with respect to $A$:

\weq{2}{\overline{O} = \overline{A}_1 + \overline{A}_2 + \cdots + \overline{A}_h,}

and consider the sums:

\weq{3}{S_i = \sum^{\mathfrak{a}_i}\frac{1}{(N\mathfrak{a}_i)^s},}

in which $s$ denotes a positive variable $> 1$, and $\mathfrak{a}_i$ runs through all the integral ideals of a class $\overline{A}_i$.

We transform the sums $S_i$ in the following way:

Since $N(\mathfrak{a}_i)$ must be one of the numbers 1, 2, 3, \dots\ we denote by $a_\nu$ the number of ideals of the class $\overline{A}_i$ whose norm $= \nu$, and obtain

\weq{4}{S_i = a_1 + \frac{a_2}{2^s} + \frac{a_3}{3^s} + \frac{a_4}{4^s} + \cdots}

where, if a number $\nu$ does not occur among the $N(\mathfrak{a}_i)$, one is to set $a_\nu = 0$.

If we denote by $Z(\nu)$ the number of ideals of the class $\overline{A}_i$ whose norm is not greater than $\nu$, then for $Z(\nu)$ we obtain from formula (6) of the preceding section a limiting value:
%CONT>>
